% 复习题 11

\begin{exblock}{复习题 11}
\item 用两种不同的方法证明: $n$ 维球坐标换元 (变换)
\[
\begin{cases}
x_{1}=r\cos\theta_{1},\\
x_{2}=r\sin\theta_{1}\cos\theta_{2},\\
x_{3}=r\sin\theta_{1}\sin\theta_{2}\cos\theta_{3},\\
\vdots\\
x_{n-1}=r\sin\theta_{1}\sin\theta_{2}\sin\theta_{3}\cdots\sin\theta_{n-2}\sin\theta_{n-1}\cos\theta_{n-1},\\
x_{n}=r\sin\theta_{1}\sin\theta_{2}\sin\theta_{3}\cdots\sin\theta_{n-2}\sin\theta_{n-1}
\end{cases}
\]
的 Jacobi 行列式为
\[
\frac{\partial(x_{1},x_{2},\cdots,x_{n})}{\partial(r,\theta_{1},\cdots,\theta_{n-1})}=r^{n-1}\sin^{n-2}\theta_{1}\sin^{n-3}\theta_{2}\cdots\sin\theta_{n-2};
\]
体积元为
\[
\mathrm{d}v=r^{n-1}\sin^{n-2}\theta_{1}\sin^{n-3}\theta_{2}\cdots\sin\theta_{n-2}\,\mathrm{d}r\mathrm{d}\theta_{1}\cdots\mathrm{d}\theta_{n-2}.
\]

  \begin{enumerate}[label=(\arabic*),leftmargin=2.6em,itemsep=0.15\baselineskip,topsep=0.2\baselineskip]

  \item 10.5 节中的证法;

  \item 记 $\boldsymbol{x}(r,\theta_{1},\cdots,\theta_{n-1})=(x_{1}(r,\theta_{1},\cdots,\theta_{n-1}),x_{2}(r,\theta_{1},\cdots,\theta_{n-1}),\cdots,x_{n}(r,\theta_{1},\cdots,\theta_{n-1}))$, 由 $\boldsymbol{x}'_{r},\boldsymbol{x}'_{\theta_{1}},\cdots,\boldsymbol{x}'_{\theta_{n-1}}$ 彼此正交推得.

  \end{enumerate}

\item 用两种不同的方法证明: $n-1$ 维球面坐标
\[
\begin{cases}
x_{1}=R\cos\theta_{1},\\
x_{2}=R\sin\theta_{1}\cos\theta_{2},\\
x_{3}=R\sin\theta_{1}\sin\theta_{2}\cos\theta_{3},\\
\vdots\\
x_{n-1}=R\sin\theta_{1}\sin\theta_{2}\sin\theta_{3}\cdots\sin\theta_{n-2}\sin\theta_{n-1}\cos\theta_{n-1},\\
x_{n}=R\sin\theta_{1}\sin\theta_{2}\sin\theta_{3}\cdots\sin\theta_{n-2}\sin\theta_{n-1}\sin\theta_{n-1}
\end{cases}
\]
的面积元为
\[
\mathrm{d}\sigma=R^{n-1}\sin^{n-2}\theta_{1}\sin^{n-3}\theta_{2}\cdots\sin\theta_{n-2}\,\mathrm{d}\theta_{1}\cdots\mathrm{d}\theta_{n-1}.
\]

\item 设 $\mathbb{R}^{n}$ 中 $n-1$ 维超曲面 $\Sigma$ 的参数表示为
\[
\boldsymbol{x}(x_{1},x_{2},\cdots,x_{n-1})=(x_{1},x_{2},\cdots,x_{n-1},x_{n}(x_{1},x_{2},\cdots,x_{n-1})),
\]
$(x_{1},x_{2},\cdots,x_{n-1})\in\Delta$ (开区域), 其中 $x_{n}(x_{1},x_{2},\cdots,x_{n-1})$ 有连续偏导数. 证明: $\Sigma$ 的面积元为
\[
\mathrm{d}\sigma=\sqrt{1+\left(\frac{\partial x_{n}}{\partial x_{1}}\right)^{2}+\left(\frac{\partial x_{n}}{\partial x_{2}}\right)^{2}+\cdots+\left(\frac{\partial x_{n}}{\partial x_{n-1}}\right)^{2}}\,\mathrm{d}x_{1}\cdots\mathrm{d}x_{n-1}.
\]

\item 设 $\boldsymbol{p}=(x,y,z)$, $p=\|\boldsymbol{p}\|$, $\boldsymbol{F}=\frac{\boldsymbol{p}}{p^{3}}$, $\Sigma$ 为 $\mathbb{R}^{3}$ 中 $C^{1}$ 正则的封闭曲面. 证明: 当原点在曲面 $\Sigma$ 的外, 上, 内时, 分别有
\[
\oint_{\Sigma}\boldsymbol{F}\cdot\boldsymbol{n}_{0}\,\mathrm{d}\sigma=0,\ 2\pi,\ 4\pi.
\]
其中 $\boldsymbol{n}_{0}$ 为方向朝外的连续的单位法向量场.

\item 证明: 第二型曲线积分
\[
\int_{\overrightarrow{C}}y\,\mathrm{d}x+z\,\mathrm{d}y+x\,\mathrm{d}z=\frac{ac}{2}-\frac{\pi b}{4\sqrt{2}}(a+c).
\]
其中 $\overrightarrow{C}$ 是曲线
\[
\frac{x^{2}}{a^{2}}+\frac{y^{2}}{b^{2}}+\frac{z^{2}}{c^{2}}=1,\quad\frac{x}{a}+\frac{z}{c}=1,\quad x\geqslant0,\ y\geqslant0,\ z\geqslant0
\]
($a>0,b>0,c>0$ 为常数)从点 $(a,0,0)$ 到 $(0,0,c)$ 的一段.

\item 证明: 第二型曲线积分
\[
\oint_{\overrightarrow{C}}\frac{-y\,\mathrm{d}x+x\,\mathrm{d}y}{\bigl[(\alpha x+\beta y)^{2}+(\gamma x+\delta y)^{2}\bigr]^{\alpha}}=\frac{2\pi}{|\alpha\delta-\beta\gamma|},\quad\alpha\delta-\beta\gamma\neq0,
\]
其中 $\overrightarrow{C}$ 为椭圆 $(\alpha x+\beta y)^{2}+(\gamma x+\delta y)^{2}=1$, 取逆时针方向.

\item 设 $C$ 为平面上一条自身不相交的 $C^{1}$ 曲线. 其起点为 $(1,0)$, 终点为 $(0,2)$. 除起, 终点外, $C$ 全部落在第一象限内. 证明: 第一型曲线积分
\[
\int_{C}\frac{\partial\ln r}{\partial\boldsymbol{n}_{0}}\,\mathrm{d}s=-\frac{\pi}{2},
\]
其中 $\boldsymbol{n}_{0}$ 为 $C$ 上的连续的单位法向量场, 方向指向原点所在的一侧; $r$ 为 $C$ 上动点到原点的距离, $\mathrm{d}s$ 为 $C$ 的弧长元.

\item 设 $y>0$, 求
\[
P\,\mathrm{d}x+Q\,\mathrm{d}y=\frac{y\,\mathrm{d}x-x\,\mathrm{d}y}{3x^{2}-2xy+3y^{2}}
\]
的原函数.

\item 证明: 第一型曲面积分
\[
F(t)=\iint_{x+y+z=t}f(x,y,z)\,\mathrm{d}\sigma=\begin{cases}
\frac{\pi}{18}(3-t^{2})^{2}, & |t|\leqslant\sqrt{3},\\
0, & |t|>\sqrt{3},
\end{cases}
\]
其中
\[
f(x,y,z)=\begin{cases}
1-x^{2}-y^{2}-z^{2}, & x^{2}+y^{2}+z^{2}\leqslant1,\\
0, & x^{2}+y^{2}+z^{2}>1.
\end{cases}
\]

\item 设 $f(t)$ 在 $|t|\leqslant\sqrt{a^{2}+b^{2}+c^{2}}$ 内连续. 证明: Poisson 公式
\[
\int_{0}^{2\pi}\mathrm{d}\varphi\int_{0}^{\pi}f(a\sin\theta\cos\varphi+b\sin\theta\sin\varphi+c\cos\theta)\sin\theta\,\mathrm{d}\theta=2\pi\int_{-1}^{1}f(kz)\,\mathrm{d}z,
\]
其中 $k=\sqrt{a^{2}+b^{2}+c^{2}}\neq0$.

\item 设 $f(x)$ 在 $|x|\leqslant\sqrt{a^{2}+b^{2}+c^{2}}$ ($a^{2}+b^{2}+c^{2}\neq0$) 上连续. 证明:
\[
\iiint_{\Omega}f\left(\frac{ax+by+cz}{\sqrt{x^{2}+y^{2}+z^{2}}}\right)\mathrm{d}x\mathrm{d}y\mathrm{d}z=\frac{2\pi}{3}\int_{-1}^{1}f(u\sqrt{a^{2}+b^{2}+c^{2}})\,\mathrm{d}u,
\]
其中 $\Omega$ 为单位球体 $x^{2}+y^{2}+z^{2}\leqslant1$.

\item 设 $f(t)$ 在 $|t|\leqslant1$ 内连续可导, $f(-1)=f(1)=0$, $M=\max_{-1\leqslant t\leqslant1}\{f'(t)\}$, $S^{2}(1)$ 是中心在原点半径为 1 的球面. 证明:
\[
\left|\iint_{S^{2}(1)}f(mx+ny+pz)\,\mathrm{d}\sigma\right|\leqslant2\pi M,
\]
其中 $m^{2}+n^{2}+p^{2}=1$, $m,n,p$ 为常数.

\item 设 $\Sigma$ 为
\[
z-c=\sqrt{R^{2}-(x-a)^{2}-(y-b)^{2}}
\]
的上侧. 证明: 第二型曲面积分
\[
\iint_{\Sigma}x^{2}\mathrm{d}y\wedge\mathrm{d}z+y^{2}\mathrm{d}z\wedge\mathrm{d}x+(x-a)yz\,\mathrm{d}x\wedge\mathrm{d}y=\frac{4}{3}\pi(a+b)R^{3}.
\]

\item 设 $\Sigma$ 为曲面 $|x-y+z|+|y-z+x|+|z-x+y|=1$ 的外表面. 证明: 第二型曲面积分
\[
\iint_{\Sigma}(x+y-z)\mathrm{d}y\wedge\mathrm{d}z+\bigl[2y+\sin(z+x)\bigr]\mathrm{d}z\wedge\mathrm{d}x+(3z+\mathrm{e}^{x+y})\mathrm{d}x\wedge\mathrm{d}y=2.
\]

\item 设 $\overrightarrow{\Sigma}$ 为 $1-\frac{z}{7}=\frac{(x-2)^{2}}{25}+\frac{(y-1)^{2}}{16}$ ($z\geqslant0$) 的上侧. 证明: 第二型曲面积分
\[
\iint_{\Sigma}\frac{x\,\mathrm{d}y\wedge\mathrm{d}z+y\,\mathrm{d}z\wedge\mathrm{d}x+z\,\mathrm{d}x\wedge\mathrm{d}y}{\sqrt{(x^{2}+y^{2}+z^{2})^{3}}}=2\pi.
\]
\end{exblock}
